%=====================================================================
\chapter{Sub-queries, \texorpdfstring{\sqlkw{EXISTS}}{EXISTS} and Set Operations}\label{ch:subqueries}
%=====================================================================
\locator{4}{Part IV --- Asking Questions}

\begin{outcomes}
By the end of this chapter you will be able to:
\begin{tight}
  \item \textbf{Write} a sub-query in the \sqlkw{WHERE}, \sqlkw{FROM} and
        \sqlkw{SELECT} clauses, and \textbf{say} what each one must return.
  \item \textbf{Distinguish} a correlated from an uncorrelated sub-query by
        reading it.
  \item \textbf{Write} a semi-join and an anti-join in several forms, and
        \textbf{choose} one that is correct whatever the nulls.
  \item \textbf{Explain} what makes a correlated sub-query expensive, and
        \textbf{show} that it is not the correlation.
  \item \textbf{Use} a common table expression to give a sub-query a name.
\end{tight}
\end{outcomes}

\begin{prereq}
\chref{select}, and especially its \sqlkw{NOT IN} alert --- this chapter
measures the four alternatives. \chref{joins} for the anti-join idiom. You
need \texttt{code/m15\_subqueries.py}.
\end{prereq}

\begin{keyterms}
sub-query $\bullet$ nested query $\bullet$ scalar sub-query $\bullet$
correlated sub-query $\bullet$ uncorrelated sub-query $\bullet$ semi-join
$\bullet$ anti-join $\bullet$ \sqlkw{IN} $\bullet$ \sqlkw{EXISTS} $\bullet$
\sqlkw{ANY} $\bullet$ \sqlkw{ALL} $\bullet$ derived table $\bullet$ common
table expression $\bullet$ \sqlkw{WITH} $\bullet$ \sqlkw{UNION} $\bullet$
\sqlkw{UNION ALL} $\bullet$ \sqlkw{INTERSECT} $\bullet$ \sqlkw{EXCEPT}
\end{keyterms}

\section{A Query Inside a Query}

\chref{algebra}'s closure property --- every result is a relation --- is what
makes this chapter possible. A \sqlkw{SELECT} produces a relation, so a
\sqlkw{SELECT} can appear wherever a relation, a list or a single value is
expected.

\begin{definitionbox}[Three places a sub-query can go]
\textbf{In \sqlkw{WHERE}}, after \sqlkw{IN} or \sqlkw{EXISTS}: it supplies a
list of values, or the fact that some row exists.

\smallskip
\textbf{In \sqlkw{FROM}}: it supplies a whole relation, called a
\term{derived table}, which must be given an alias.

\smallskip
\textbf{In the \sqlkw{SELECT} list}, or in a comparison: it must be a
\term{scalar sub-query}, returning exactly one column and at most one row. If
it returns no rows its value is null; if it returns two, SQLite silently uses
the first --- other systems raise an error.
\end{definitionbox}

\begin{definitionbox}[Correlated and uncorrelated]
An \term{uncorrelated} sub-query mentions only its own tables. It could be run
on its own, and in principle it is run once.

\smallskip
A \term{correlated} sub-query mentions a column of the \emph{outer} query, so
its answer depends on which outer row is being considered. In principle it is
run once per outer row.

\begin{lstlisting}[language=sql]
-- uncorrelated: the inner query stands alone
WHERE roll_no IN (SELECT roll_no FROM enrolment WHERE grade = 'A')

-- correlated: s.roll_no belongs to the OUTER query
WHERE EXISTS (SELECT 1 FROM enrolment e
               WHERE e.roll_no = s.roll_no AND e.grade = 'A')
\end{lstlisting}

\smallskip
``In principle'' is doing work in both definitions. The optimiser is free to
run either one however it likes, and the measurement below shows it doing so.
\end{definitionbox}

\section{Semi-joins and Anti-joins}

\begin{definitionbox}[Semi-join, anti-join]
A \term{semi-join} returns the rows of $R$ that have \emph{at least one}
partner in $S$ --- each row once, however many partners it has. ``Students with
at least one A.''

\smallskip
An \term{anti-join} returns the rows of $R$ that have \emph{no} partner in
$S$. ``Instructors who teach nothing.'' It is \chref{algebra}'s difference,
and ``never'', ``none'' and ``no'' always mean one.
\end{definitionbox}

\begin{measured}[Seven spellings, one wrong answer, and correlation that was not the problem]
\textsf{\footnotesize \texttt{code/m15\_subqueries.py}, seeded database,
SQLite 3.51.2, Intel i5-6600, Windows 11. Best of three runs.}

\medskip
\textbf{1. A semi-join, three ways} --- students with at least one A

\begin{tabular}{@{}L{5.0cm}R{2.0cm}R{1.9cm}L{5.0cm}@{}}
\textbf{written as} & \textbf{answer} & \textbf{ms} & \\
\sqlkw{IN} (uncorrelated)   & 3{,}449 & 20.3 & \\
\sqlkw{EXISTS} (correlated) & 3{,}449 & 40.1 & twice the \sqlkw{IN} \\
\sqlkw{JOIN} + \sqlkw{DISTINCT} & 3{,}449 & 21.4 & \\
\end{tabular}

\medskip
\textbf{2. An anti-join, four ways} --- instructors who teach nothing, after
three are hired who teach nothing yet; \texttt{section.emp\_id} has 73 nulls

\begin{tabular}{@{}L{5.0cm}R{2.0cm}R{1.9cm}L{5.0cm}@{}}
\textbf{written as} & \textbf{answer} & \textbf{ms} & \textbf{right?} \\
\sqlkw{NOT IN}                 & \textbf{0} & 0.76 & \textbf{no} --- one null empties it \\
\sqlkw{NOT EXISTS}             & 3 & 2.18 & yes \\
\sqlkw{LEFT JOIN} \dots\ \sqlkw{IS NULL} & 3 & 2.52 & yes \\
\sqlkw{EXCEPT}                 & 3 & 0.44 & yes \\
\end{tabular}

\smallskip
The fastest wrong answer is no use. \sqlkw{EXCEPT} is a set operation rather
than a comparison, so the null on its right-hand side can only remove a null
on its left --- and \texttt{instructor.emp\_id}, a primary key, has none.

\medskip
\textbf{3. A correlated sub-query in the \sqlkw{SELECT} list} --- enrolments
per course

\begin{tabular}{@{}R{1.7cm}R{2.1cm}R{2.4cm}R{2.1cm}R{2.4cm}R{2.4cm}@{}}
& & \multicolumn{2}{c}{\textbf{no index on course\_code}}
& \multicolumn{2}{c}{\textbf{with the index}} \\
\textbf{students} & \textbf{enrolments} & \textbf{correlated} & \textbf{\sqlkw{GROUP BY}}
& \textbf{correlated} & \textbf{\sqlkw{GROUP BY}} \\
2{,}000 & 30{,}027  & 191.1 & 53.8  & \textbf{2.4}  & 5.4 \\
4{,}000 & 60{,}037  & 517.7 & 121.3 & \textbf{5.1}  & 12.2 \\
8{,}000 & 120{,}134 & 1{,}101.4 & 275.8 & \textbf{10.3} & 23.7 \\
\end{tabular}

\smallskip
All four agree on every count. Times in milliseconds.

\medskip
\textbf{Read part 3 carefully, because it contradicts the folklore.} Without
an index the correlated form is about four times slower than \sqlkw{GROUP BY}:
\textit{enrolment}'s primary key is \texttt{(roll\_no, course\_code, \dots)},
so it cannot be searched by \texttt{course\_code}, and each of the 48 courses
scans all 60 thousand enrolments. \textbf{Add one index and the correlated form
becomes the faster of the two} --- 5.1 against 12.2\,ms --- because it now does
48 index lookups, while \sqlkw{GROUP BY} still reads every enrolment. Both
grow linearly with the data; neither diverges from the other as the table
grows.

\smallskip
\textbf{So the cost was never the correlation.} It was whether the inner query
had an index on the column it is correlated on. The plans say so directly:
\texttt{SCAN e} inside \texttt{CORRELATED SCALAR SUBQUERY} without it,
\texttt{SEARCH e USING \dots\ idx\_enr\_course (course\_code=?)} with it.
\end{measured}

\begin{alertbox}[What to take from the measurement]
\begin{tight}
  \item \textbf{For anti-joins, use \sqlkw{NOT EXISTS}.} It is correct
        whatever the nulls, and its cost is in the same range as the
        alternatives. \sqlkw{EXCEPT} is faster when it fits but returns only
        the compared columns; \sqlkw{LEFT JOIN} \dots\ \sqlkw{IS NULL} is
        correct but easy to break with a condition in the wrong clause
        (\chref{joins}).
  \item \textbf{Do not avoid correlated sub-queries on principle.} Avoid
        unindexed ones. ``Correlated sub-queries are slow'' is a rule of thumb
        learned on databases without the right index, and on this one it was
        backwards by a factor of two.
  \item \textbf{Read the plan before rewriting a query.} \sqlkw{SCAN} inside a
        \sqlkw{CORRELATED} sub-query is the specific thing to look for; it
        means the inner query reads a whole table once per outer row.
\end{tight}
\end{alertbox}

\section{\sqlkw{ANY}, \sqlkw{ALL}, and Naming a Sub-query}

\begin{examplebox}[\sqlkw{ALL} and \sqlkw{ANY}, and the null trap again]
\texttt{x > ALL (sub-query)} is true when $x$ exceeds every value returned;
\texttt{x > ANY (sub-query)} when it exceeds at least one. SQLite does not
support either keyword, which is no loss --- both are clearer as an aggregate:

\begin{lstlisting}[language=sql]
-- "credit hours greater than every MA course"
WHERE credit_hours > (SELECT MAX(credit_hours) FROM course
                       WHERE dept_code = 'MA')
\end{lstlisting}

\smallskip
And \sqlkw{ALL} has \sqlkw{NOT IN}'s trap: \texttt{x > ALL (\dots, NULL)} is
\unk{} for every $x$, because one of its comparisons is. The
aggregate form ignores nulls, which is usually what was meant.
\end{examplebox}

\begin{definitionbox}[Common table expression]
A \term{common table expression} (CTE) names a sub-query before the main query
uses it:

\begin{lstlisting}[language=sql]
WITH a_students AS (
  SELECT DISTINCT roll_no FROM enrolment WHERE grade = 'A'
)
SELECT s.name
  FROM student s JOIN a_students a ON a.roll_no = s.roll_no
 WHERE s.programme = 'BSIT';
\end{lstlisting}

\smallskip
It computes nothing a derived table could not. Its value is that the
sub-query has a \emph{name}, can be read top to bottom, and can be referred to
more than once. \sqlkw{WITH RECURSIVE} --- used in \chref{whodoeswhat}'s
laboratory to generate 300 rows --- is the one form that does something new:
it can walk the prerequisite chain to any depth, which no fixed number of joins
can.
\end{definitionbox}

\begin{worked}[Every prerequisite of a course, to any depth]
\textbf{The question.} Before taking a course, which courses must a student
have passed --- directly, or because a prerequisite has a prerequisite?

\medskip
\textbf{Step 1 --- why a join cannot do it.} One join finds direct
prerequisites. Two joins find their prerequisites. Each extra level needs an
extra join, and nobody knows in advance how deep the chain goes.

\medskip
\textbf{Step 2 --- the recursive CTE.} Start from the direct prerequisites
(the \emph{anchor}), then repeatedly add the prerequisites of whatever is
already found (the \emph{recursive step}), until a step adds nothing:

\begin{lstlisting}[language=sql]
WITH RECURSIVE needs(code) AS (
  SELECT requires FROM prerequisite
   WHERE course_code = 'CS-305'                 -- the anchor
  UNION
  SELECT p.requires FROM prerequisite p
    JOIN needs n ON p.course_code = n.code      -- the step
)
SELECT code FROM needs;
\end{lstlisting}

\medskip
\textbf{Step 3 --- the answer, on the seeded database.} \texttt{CS-305}
directly requires two courses, \texttt{CS-201} and \texttt{CS-258}. The query
returns \textbf{four}: those two, plus \texttt{CS-150} (required by
\texttt{CS-258}) and \texttt{CS-101} (required by \texttt{CS-150}). The chain
is three levels deep --- 305, 258, 150, 101 --- and a query with one join
would have found half of it.

\medskip
\textbf{Step 4 --- why \sqlkw{UNION} and not \sqlkw{UNION ALL}.} \sqlkw{UNION}
discards rows already found, so a course reachable by two paths is listed
once --- and, more importantly, a cycle in the data (A requires B requires A)
terminates instead of running for ever. \textbf{\chref{mapping}'s
\sqlkw{CHECK} stops a course requiring itself, but not a two-step cycle.} Only
\sqlkw{UNION} saves you from that one.
\end{worked}

\dsfig{%
\begin{tikzpicture}[
  crs/.style={rounded corners=3pt, draw=#1, line width=1pt, fill=#1!8,
              text=#1!60!black, font=\scriptsize\bfseries\ttfamily,
              minimum width=1.6cm, minimum height=0.62cm},
  req/.style={-{Stealth[length=2.2mm]}, neutral!70, line width=0.9pt},
  lvl/.style={font=\scriptsize\itshape, text=neutral},
  br/.style={draw=#1, line width=0.9pt,
             decorate, decoration={brace, amplitude=5pt, mirror}}
]
  \node[crs=primary]   (a) at (0,0)    {CS-305};
  \node[crs=greenhl]   (b) at (3.0,0.75) {CS-201};
  \node[crs=greenhl]   (c) at (3.0,-0.75) {CS-258};
  \node[crs=goldhl]    (d) at (6.0,-0.75) {CS-150};
  \node[crs=accent]    (e) at (9.0,-0.75) {CS-101};
  \draw[req] (a) -- (b);
  \draw[req] (a) -- (c);
  \draw[req] (c) -- (d);
  \draw[req] (d) -- (e);
  \node[lvl, anchor=south] at (3.0,1.25) {level 1};
  \node[lvl, anchor=south] at (6.0,-0.15) {level 2};
  \node[lvl, anchor=south] at (9.0,-0.15) {level 3};
  \draw[br=greenhl] (2.1,-1.35) -- (3.9,-1.35);
  \node[lvl, text=greenhl!60!black, anchor=north] at (3.0,-1.6)
       {one join finds these two};
  \draw[br=accent] (2.1,-2.25) -- (9.9,-2.25);
  \node[lvl, text=accent, anchor=north] at (6.0,-2.5)
       {\sqlkw{WITH RECURSIVE} finds all four, however deep};
\end{tikzpicture}}%
{The prerequisites of \texttt{CS-305} on the seeded database. An arrow means
``requires''. A fixed query reaches a fixed depth; the recursive one keeps
going until a step adds nothing new.}{prereqchain}

\begin{pitfall}
\begin{tight}
  \item \textbf{\sqlkw{NOT IN} against a nullable column.} Measured: 0 instead
        of 3.
  \item \textbf{Assuming correlation is the cost.} The missing index is.
        Measured: with the index, the correlated form was the faster.
  \item \textbf{A scalar sub-query that returns two rows.} SQLite quietly uses
        one; other systems refuse.
  \item \textbf{A derived table with no alias.} Required.
  \item \textbf{\sqlkw{UNION} where \sqlkw{UNION ALL} was meant}, or the
        reverse. \sqlkw{UNION} removes duplicates and costs a sort; in a
        recursive CTE it is also what stops a cycle.
  \item \textbf{\sqlkw{x > ALL} over a list with a null.} \unk{}
        for every $x$. Use \sqlkw{MAX}.
\end{tight}
\end{pitfall}

\begin{lab}[Write one question seven ways, and find the one that lies]
\begin{lstlisting}[language=bash]
# --- 1. The seeded database ----------------------------------------
python code/schema/seed.py university.db
sqlite3 university.db
\end{lstlisting}
\begin{lstlisting}[language=sql]
-- --- 2. A semi-join three ways ---------------------------------
.headers on
.mode box
.timer on
SELECT COUNT(*) FROM student WHERE roll_no IN
  (SELECT roll_no FROM enrolment WHERE grade = 'A');
SELECT COUNT(*) FROM student s WHERE EXISTS
  (SELECT 1 FROM enrolment e WHERE e.roll_no = s.roll_no
                               AND e.grade = 'A');
SELECT COUNT(DISTINCT s.roll_no) FROM student s
  JOIN enrolment e ON e.roll_no = s.roll_no WHERE e.grade = 'A';
-- Same answer three times? Record the three timings.

-- --- 3. Hire three people, then an anti-join four ways -----------
PRAGMA foreign_keys = ON;
INSERT INTO instructor VALUES
  ('E-9001','New Hire One','Lecturer','CS'),
  ('E-9002','New Hire Two','Lecturer','IT'),
  ('E-9003','New Hire Three','Lecturer','SE');
SELECT COUNT(*) FROM instructor
 WHERE emp_id NOT IN (SELECT emp_id FROM section);
SELECT COUNT(*) FROM instructor i WHERE NOT EXISTS
  (SELECT 1 FROM section s WHERE s.emp_id = i.emp_id);
SELECT COUNT(*) FROM instructor i LEFT JOIN section s
    ON s.emp_id = i.emp_id WHERE s.emp_id IS NULL;
SELECT COUNT(*) FROM (SELECT emp_id FROM instructor
                      EXCEPT SELECT emp_id FROM section);
-- Three say 3. One says 0. Which, and why?

-- --- 4. A correlated sub-query, and its plan ---------------------
EXPLAIN QUERY PLAN
SELECT c.course_code, (SELECT COUNT(*) FROM enrolment e
                        WHERE e.course_code = c.course_code)
  FROM course c;
-- Find SCAN inside CORRELATED. Then time it:
SELECT c.course_code, (SELECT COUNT(*) FROM enrolment e
                        WHERE e.course_code = c.course_code)
  FROM course c;

-- --- 5. Give the inner query an index, and look again -------------
CREATE INDEX idx_enr_course ON enrolment(course_code);
EXPLAIN QUERY PLAN
SELECT c.course_code, (SELECT COUNT(*) FROM enrolment e
                        WHERE e.course_code = c.course_code)
  FROM course c;
-- SCAN has become SEARCH. Time it again and compare.

-- --- 6. Every prerequisite, to any depth ------------------------
SELECT requires FROM prerequisite WHERE course_code = 'CS-305';
-- The direct prerequisites: two. Now every one, to any depth:
WITH RECURSIVE needs(code) AS (
  SELECT requires FROM prerequisite WHERE course_code = 'CS-305'
  UNION
  SELECT p.requires FROM prerequisite p JOIN needs n
      ON p.course_code = n.code)
SELECT code FROM needs;
-- Four. Draw the chain on paper and check every one.
\end{lstlisting}
\begin{lstlisting}[language=bash]
# --- 7. The measured version ----------------------------------------
python code/m15_subqueries.py
\end{lstlisting}

\smallskip
\textbf{What to notice.} Step 5 is where a rule of thumb meets a measurement.
The correlated query you were told to avoid becomes faster than the
\sqlkw{GROUP BY} you were told to prefer, because of one index and nothing
else. Rules about query \emph{shape} are usually rules about indexes in
disguise, and \chref{indexes} takes the disguise off.
\end{lab}

\begin{checkpoint}
\begin{tightnum}
  \item Classify each sub-query in the laboratory as correlated or
        uncorrelated, and for each correlated one name the outer column it
        refers to.
  \item Of the four anti-join forms, one returned 0 and three returned 3. Say
        which, explain why in terms of three-valued logic, and say why
        \sqlkw{EXCEPT} was not affected by the same nulls.
  \item Without an index the correlated sub-query was four times slower than
        \sqlkw{GROUP BY}; with one it was twice as fast. Explain both results
        in terms of what each query reads, and say what this tells you about
        the advice ``avoid correlated sub-queries''.
\end{tightnum}
\end{checkpoint}

\begin{chaptersummary}
\begin{tight}
  \item A sub-query can go in \sqlkw{WHERE} (a list, or existence),
        \sqlkw{FROM} (a derived table, which needs an alias) or the
        \sqlkw{SELECT} list (a scalar: one column, at most one row).
  \item A \textbf{correlated} sub-query mentions an outer column and, in
        principle, runs per outer row. The optimiser decides what actually
        happens.
  \item A \textbf{semi-join} keeps rows with at least one partner; an
        \textbf{anti-join} keeps rows with none.
  \item \textbf{Measured: the three semi-join forms agreed} (3{,}449);
        \sqlkw{EXISTS} took twice as long as \sqlkw{IN} here.
  \item \textbf{Measured: \sqlkw{NOT IN} returned 0} where \sqlkw{NOT
        EXISTS}, \sqlkw{LEFT JOIN} \dots\ \sqlkw{IS NULL} and \sqlkw{EXCEPT}
        all returned 3. Use \sqlkw{NOT EXISTS}.
  \item \textbf{Measured: without an index on the correlation column, the
        correlated form was 4$\times$ slower than \sqlkw{GROUP BY}; with one
        it was 2$\times$ faster.} The cost was never the correlation.
  \item \sqlkw{SCAN} inside a \sqlkw{CORRELATED} sub-query in a plan means a
        whole table read per outer row. Look for it.
  \item A \textbf{CTE} names a sub-query. \sqlkw{WITH RECURSIVE} walks a chain
        to any depth, and \sqlkw{UNION} in it is what stops a cycle.
  \item \sqlkw{ALL} over a list containing a null is \unk{}; use an
        aggregate.
\end{tight}
\end{chaptersummary}

\begin{reviewq}
  \item \textbf{(MCQ)} A sub-query that mentions a column of the outer query
        is: (a)~scalar; (b)~derived; (c)~correlated; (d)~recursive.
  \item \textbf{(MCQ)} Which anti-join form can return no rows because of a
        single null? (a)~\sqlkw{NOT EXISTS}; (b)~\sqlkw{NOT IN};
        (c)~\sqlkw{EXCEPT}; (d)~\sqlkw{LEFT JOIN} \dots\ \sqlkw{IS NULL}.
  \item \textbf{(Short)} What must a scalar sub-query return, and what does
        SQLite do if it returns more?
  \item \textbf{(Short)} Why does a recursive CTE over \textit{prerequisite}
        need \sqlkw{UNION} rather than \sqlkw{UNION ALL}?
  \item \textbf{(Short)} State in one sentence what decides whether a
        correlated sub-query is expensive.
  \item \textbf{(Applied)} Write ``courses with no enrolments in FALL2026''
        as \sqlkw{NOT EXISTS}, as \sqlkw{LEFT JOIN} \dots\ \sqlkw{IS NULL} and
        as \sqlkw{EXCEPT}. Run all three on the seeded database (insert one
        new course first so the answer is not empty), confirm they agree, and
        read each plan to say which index each one uses.
  \item \textbf{(Applied)} A student's prerequisite chain is to be checked at
        enrolment time. Write the recursive query that lists every course a
        student must have passed before taking a given one, then extend it to
        list only those the student has \emph{not} passed. Say what happens if
        somebody inserts a two-step cycle, and how your query behaves.
\end{reviewq}
